\documentclass[11pt,a4paper]{article}

\usepackage[utf8]{inputenc}
\usepackage[T1]{fontenc}
\usepackage[spanish]{babel}
\usepackage{lmodern}
\usepackage[margin=2.5cm]{geometry}
\usepackage{tabularx,array,fancyhdr}
\usepackage[table]{xcolor}
\usepackage[hidelinks]{hyperref}

\setlength{\headheight}{14pt}
\setlength{\parindent}{0pt}
\setlength{\parskip}{6pt}
\setlength{\emergencystretch}{1.5em}
\renewcommand{\arraystretch}{1.23}

\pagestyle{fancy}
\fancyhf{}
\fancyhead[L]{\small UTEQ \textperiodcentered{} Ingeniería de Requerimientos}
\fancyhead[R]{\small GA-PRAC-01}
\fancyfoot[L]{\small Grupo CEV}
\fancyfoot[R]{\small Página \thepage}

\addto\captionsspanish{\renewcommand{\refname}{Referencias}}
\usepackage[backend=biber,style=ieee]{biblatex}
\addbibresource{referencias.bib}
\begin{document}

\begin{center}
    {\large\textbf{UNIVERSIDAD TÉCNICA ESTATAL DE QUEVEDO}}\par

    Facultad de Ciencias de la Computación\\
    Carrera de Software\par

    \vspace{0.6cm}

    {\LARGE\bfseries
        Comparación entre modelos de proceso\\
        de requisitos en escenarios reales\par
    }

    \vspace{0.2cm}

    {\large GA-PRAC-01}\\
    Actividad formativa grupal
\end{center}

\vspace{0.4cm}

\section{Datos informativos}

\begin{tabularx}{\textwidth}{|l|X|}
    \hline
    \rowcolor{gray!15}
    \textbf{Campo} & \textbf{Detalle} \\ \hline
    Integrantes \textperiodcentered{} CEV &
    Contreras Chavez Kevin German\newline
    Espinoza Pivaque Angelo Maximiliano\newline
    Vaca Romero David Octavio \\ \hline
    Asignatura & Ingeniería de Requerimientos (ISR-401) \\ \hline
    Nivel y paralelo & Cuarto nivel, paralelo A \\ \hline
    Docente & Gleiston Cicerón Guerrero Ulloa, PhD \\ \hline
    Código de la actividad & GA-PRAC-01 --- FORM \\ \hline
    Fecha límite de entrega & Viernes 13 de noviembre de 2026, 00:00 \\ \hline
\end{tabularx}

\section{Caracterización de los modelos de proceso de requisitos}

La comparación considera cuándo se definen los requisitos, cómo se
incorporan los cambios y en qué momentos participan las partes
interesadas. La caracterización se basa en la sección 2 y la Tabla 2
del material complementario de la práctica \cite{materialGA}.

\subsection{Modelo secuencial o en cascada}

El proyecto se organiza en fases: definición de requisitos, diseño,
implementación y pruebas. Al inicio, el equipo obtiene, analiza y
documenta los requisitos con las partes interesadas. Una vez aprobada,
la especificación sirve de base para el diseño y las pruebas. Si luego
se solicita un cambio, se revisa su impacto antes de actualizar lo
acordado. El modelo supone que los requisitos principales se pueden
definir temprano y que el cliente puede participar en su aprobación.
Conviene cuando el alcance es estable y se requieren documentos
aprobados para avanzar entre fases \cite{materialGA}.

\subsection{Modelo iterativo e incremental}

El equipo parte de un conjunto inicial de requisitos y lo refina en
ciclos. En cada ciclo se selecciona una parte del alcance, se precisan
sus requisitos y se construye un avance que las partes interesadas
pueden revisar. La iteración permite corregir o precisar lo que ya se
había definido, mientras que el incremento añade una parte funcional
al producto. Este modelo conviene cuando el objetivo general está
claro, quedan detalles por resolver y el cliente puede revisar avances
al cierre de cada ciclo \cite{materialGA}.

\subsection{Modelo ágil}

Los requisitos se gestionan en una lista priorizada que se actualiza
a medida que el equipo aprende sobre el producto. Antes de cada ciclo
corto, se detallan los elementos próximos a desarrollarse y se
acuerdan criterios de aceptación. Al revisar el incremento, las
partes interesadas dan retroalimentación y el representante del
cliente puede ajustar las prioridades. Este modelo conviene cuando
se esperan cambios frecuentes y hay una persona con disponibilidad
y autoridad para aclarar necesidades y decidir prioridades
\cite{materialGA}.

\printbibliography[title={Referencias}]

\section*{Declaración de uso de inteligencia artificial}

Se utilizó ChatGPT como apoyo para organizar las ideas de la sección
sobre los modelos de proceso de requisitos. La revisión se centró en
ordenar los apartados, corregir faltas de ortografía y puntuación, y
mejorar los conectores entre párrafos para dar continuidad a la explicación.

\end{document}